\documentclass[11pt,a4paper]{article}
\usepackage[utf8]{inputenc}\usepackage[T1]{fontenc}\usepackage[italian,english]{babel}
\usepackage{amsmath,amssymb,amsthm}\usepackage[margin=2.2cm]{geometry}\usepackage{longtable,booktabs,array,calc}
\usepackage{listings,xcolor}\usepackage{hyperref}\usepackage{url}
\providecommand{\tightlist}{\setlength{\itemsep}{0pt}\setlength{\parskip}{0pt}}
\lstset{basicstyle=\ttfamily\scriptsize,breaklines=true,frame=single,captionpos=b,columns=fullflexible,keepspaces=true,inputencoding=utf8,extendedchars=true}
\title{Proof Pursuit --- Problem 1, part 2\\ \large prove, decisioni degli agenti, codice verificato, letteratura}
\author{Team Proof Pursuit (Researcher: T. Tumini; Referee: gabundos; Creative: M. Renzi; gio3r3gio)}
\date{\today}
\begin{document}\maketitle\tableofcontents
\section*{Come leggere questo rapporto}
Per ogni cella: la richiesta ufficiale, lo stato dichiarato dal team (mai ``risolto'' senza revisione umana), la consegna
con la prova, poi la traccia delle decisioni degli agenti (Researcher: famiglia e motivo; Referee: verdetto ed errore
fatale; Creative: quando e perché è intervenuto) e il codice su cui il tentativo poggia con l'esito della riesecuzione
fidata. DIMOSTRATO, CITATO, VERIFICATO SU INTERVALLO FINITO e CONGETTURA restano distinti. L'architettura del sistema
è descritta nel README del repository.
\section{Problem 1 — Angles between lines (Fejes Toth)}
\subsection{Cella 2 (2 punti) — stato: revisionato}
\subparagraph{Parte 2 --- An orthogonality
lemma}\label{parte-2-an-orthogonality-lemma}

\textbf{Punteggio:} 2 points · \textbf{Valutazione:} Judged

A statement about a chain of vectors, in which each vector may fail to
be orthogonal only to its immediate neighbours in the list. Let
\(m \ge 2\) and let \(x_1, \ldots, x_m\) be unit vectors in
\(\mathbb{R}^{m-1}\) with \(\langle x_i, x_j \rangle = 0\) whenever
\(|i - j| \ge 2\). Prove that \[
\sum_{i=1}^{m-1} \theta(x_i, x_{i+1}) \;\le\; (m-2) \cdot \frac{\pi}{2},
\] where \(\theta(x,y) = \arccos\,\lvert \langle x, y \rangle \rvert\).

\textbf{Enunciato protetto dato agli agenti (state.json).} \texttt{Let m \textgreater{}= 2 and x\_1..x\_m unit vectors in R\textasciicircum{}\{m-1\} with \textless{}x\_i,x\_j\textgreater{} = 0 whenever |i-j| \textgreater{}= 2. Prove sum\_\{i=1\}\textasciicircum{}\{m-1\} theta(x\_i,x\_\{i+1\}) \textless{}= (m-2) pi/2, where theta(x,y) = arccos|\textless{}x,y\textgreater{}|.}

\textbf{Evidenza registrata in STATUS.md.} prova completa per induzione con proiezione (argomento F); Referee READY\_FOR\_HUMAN (A PASS, B PASS); approvata; submission/parte-2.md

\subsubsection{Consegna (prova)}
\paragraph{Problema 1 --- Parte 2 --- bozza di
consegna}\label{problema-1-parte-2-bozza-di-consegna}

\textbf{Stato dichiarato: RISOLTA (approvata).} Prova completa scritta
dal team; Referee automatico READY\_FOR\_HUMAN (giudice A matematica
PASS, giudice B evidenze PASS); approvazione umana registrata in
\texttt{runs/p1\_c2/approval.json}.

\subparagraph{1. Risultato e ambito}\label{risultato-e-ambito}

Lemma di ortogonalità per una catena di versori: l'enunciato della parte
è dimostrato integralmente, nella generalità richiesta (vedi la prova
per le ipotesi esatte). Nessun calcolo al computer è necessario.

\subparagraph{2. Dimostrazione}\label{dimostrazione}

\paragraph{Cell 2: the orthogonality lemma for a chain of unit
vectors}\label{cell-2-the-orthogonality-lemma-for-a-chain-of-unit-vectors}

\textbf{Target.} Let \(m\ge2\) and let \(x_1,\dots,x_m\) be unit vectors
in \(\mathbb R^{m-1}\) with \(\langle x_i,x_j\rangle=0\) whenever
\(|i-j|\ge2\). Then
\(\sum_{i=1}^{m-1}\theta(x_i,x_{i+1})\le(m-2)\frac{\pi}{2}\), where
\(\theta(x,y)=\arccos|\langle x,y\rangle|\in[0,\pi/2]\).

Write \(\theta_i=\theta(x_i,x_{i+1})\) for \(1\le i\le m-1\). We say
that a list of unit vectors \(y_1,\dots,y_r\) in a Euclidean space \(V\)
is a \emph{chain} if \(\langle y_i,y_j\rangle=0\) whenever
\(|i-j|\ge2\). The statement to prove is: (P\(_m\)) every chain of \(m\)
unit vectors in a Euclidean space of dimension \(m-1\) satisfies
\(\sum_{i=1}^{m-1}\theta(y_i,y_{i+1})\le(m-2)\pi/2\). (Stating it for an
arbitrary Euclidean space of dimension \(m-1\) instead of
\(\mathbb R^{m-1}\) is harmless: choosing an orthonormal basis gives a
linear isometry onto \(\mathbb R^{m-1}\), which preserves inner
products, hence the chain condition and all the angles.)

\subparagraph{Two elementary facts}\label{two-elementary-facts}

\textbf{Fact A.} For \(\beta,\psi\in[0,\pi/2]\):
\(\arccos(\cos\beta\cos\psi)\le\beta+\psi\). \emph{Proof.}
\(\cos(\beta+\psi)=\cos\beta\cos\psi-\sin\beta\sin\psi\le\cos\beta\cos\psi\)
because \(\sin\beta,\sin\psi\ge0\). Both \(\cos(\beta+\psi)\) and
\(\cos\beta\cos\psi\) lie in \([-1,1]\) and \(\arccos\) is decreasing on
\([-1,1]\), so
\(\arccos(\cos\beta\cos\psi)\le\arccos(\cos(\beta+\psi))=\beta+\psi\),
the last equality because \(\beta+\psi\in[0,\pi]\). \(\square\)

\textbf{Fact B.} If \(V\) is a Euclidean space of dimension \(n\ge1\)
and \(v_1,\dots,v_{n-1}\in V\) (possibly fewer vectors, possibly none),
there is a unit vector \(y\in V\) orthogonal to all of them.
\emph{Proof.} The orthogonal complement of
\(\mathrm{span}(v_1,\dots,v_{n-1})\) in \(V\) has dimension
\(\ge n-(n-1)=1\), so it contains a nonzero vector; normalise it.
\(\square\)

\subparagraph{\texorpdfstring{Induction on
\(m\)}{Induction on m}}\label{induction-on-m}

\textbf{Base \(m=2\).} \(x_1,x_2\) are unit vectors in a
\(1\)-dimensional space, so \(x_2=\pm x_1\),
\(|\langle x_1,x_2\rangle|=1\) and
\(\theta_1=\arccos 1=0\le0=(m-2)\pi/2\).

\textbf{Step.} Let \(m\ge3\) and assume (P\(_{m-1}\)). Let
\(x_1,\dots,x_m\) be a chain of unit vectors in a Euclidean space \(W\)
of dimension \(m-1\). Put
\(U=x_m^{\perp}=\{v\in W:\langle v,x_m\rangle=0\}\), a subspace of
dimension \(m-2\). For \(j\le m-2\) we have \(|j-m|\ge2\), hence
\(\langle x_j,x_m\rangle=0\): so \(x_1,\dots,x_{m-2}\in U\). Let
\(c=\langle x_{m-1},x_m\rangle\in[-1,1]\) and \(w=x_{m-1}-c\,x_m\). Then
\(\langle w,x_m\rangle=c-c=0\), so \(w\in U\), and
\(\|w\|^2=1-2c^2+c^2=1-c^2\).

\emph{Case 1: \(|c|<1\).} Then \(w\neq0\); let \(y=w/\|w\|\), a unit
vector of \(U\). For \(j\le m-3\) we have \(|j-(m-1)|\ge2\) and
\(|j-m|\ge2\), so
\(\langle y,x_j\rangle=(\langle x_{m-1},x_j\rangle-c\langle x_m,x_j\rangle)/\|w\|=0\).
Consequently \(x_1,\dots,x_{m-2},y\) is a chain of \(m-1\) unit vectors
in \(U\) (pairs among \(x_1,\dots,x_{m-2}\) inherit the chain condition
from the original chain; the pairs \((x_j,y)\) with \(j\le m-3\),
i.e.~\(|j-(m-1)|\ge2\) in the new indexing, were just checked). Since
\(\dim U=m-2\), (P\(_{m-1}\)) gives
\[\sum_{i=1}^{m-3}\theta_i+\psi\le(m-3)\frac{\pi}{2},\qquad\text{where }\psi=\theta(x_{m-2},y)\in[0,\pi/2].\tag{3}\]
(For \(m=3\) the sum is empty and (3) says \(\psi\le0\),
i.e.~\(\psi=0\), which is indeed what (P\(_2\)) gives for the chain
\(x_1,y\) in the \(1\)-dimensional space \(U\).) Now
\(x_{m-1}=w+c\,x_m=\|w\|\,y+c\,x_m\). Let
\(\varphi=\theta_{m-1}=\theta(x_{m-1},x_m)=\arccos|c|\), so
\(|c|=\cos\varphi\), \(\|w\|=\sqrt{1-c^2}=\sin\varphi\) and
\(\varphi\in(0,\pi/2]\). Using \(\langle x_{m-2},x_m\rangle=0\),
\[\langle x_{m-2},x_{m-1}\rangle=\|w\|\langle x_{m-2},y\rangle+c\langle x_{m-2},x_m\rangle=\sin\varphi\,\langle x_{m-2},y\rangle,\]
so
\(|\langle x_{m-2},x_{m-1}\rangle|=\sin\varphi\cos\psi=\cos\beta\cos\psi\)
with \(\beta=\pi/2-\varphi\in[0,\pi/2)\). By Fact A,
\[\theta_{m-2}=\arccos(\cos\beta\cos\psi)\le\beta+\psi=\frac{\pi}{2}-\varphi+\psi,\qquad\text{i.e.}\qquad \theta_{m-2}+\theta_{m-1}\le\frac{\pi}{2}+\psi.\]
Adding this to (3):
\(\sum_{i=1}^{m-1}\theta_i\le(m-3)\frac{\pi}{2}-\psi+\frac{\pi}{2}+\psi=(m-2)\frac{\pi}{2}\).

\emph{Case 2: \(|c|=1\), i.e.~\(x_{m-1}=\pm x_m\).} Then
\(\theta_{m-1}=\arccos1=0\), and
\(\langle x_{m-2},x_{m-1}\rangle=\pm\langle x_{m-2},x_m\rangle=0\), so
\(\theta_{m-2}=\arccos0=\pi/2\). By Fact B applied to \(V=U\) (dimension
\(m-2\ge1\)) and the \(m-3\) vectors \(x_1,\dots,x_{m-3}\), there is a
unit vector \(y\in U\) with \(\langle y,x_j\rangle=0\) for all
\(j\le m-3\). Then \(x_1,\dots,x_{m-2},y\) is a chain of \(m-1\) unit
vectors in \(U\) (same verification as in Case 1), and (P\(_{m-1}\))
gives
\(\sum_{i=1}^{m-3}\theta_i+\theta(x_{m-2},y)\le(m-3)\frac{\pi}{2}\),
hence \(\sum_{i=1}^{m-3}\theta_i\le(m-3)\frac{\pi}{2}\) because
\(\theta(x_{m-2},y)\ge0\). Therefore
\[\sum_{i=1}^{m-1}\theta_i=\sum_{i=1}^{m-3}\theta_i+\frac{\pi}{2}+0\le(m-2)\frac{\pi}{2}.\]

In both cases (P\(_m\)) holds, which completes the induction.
\(\blacksquare\)

\emph{Remark (not needed).} For \(m=3\) the argument shows
\(\theta_1+\theta_2=\pi/2\) always; for general \(m\) equality is
attained e.g.~by \(x_i=e_i\) (\(i\le m-1\)), \(x_m=e_{m-1}\), where all
consecutive angles are \(\pi/2\) except the last, which is \(0\).

\subparagraph{3. Verifica: istruzioni, dipendenze,
tempi}\label{verifica-istruzioni-dipendenze-tempi}

Prova a mano, nessuna dipendenza. Verifica meccanica non necessaria; il
Referee automatico (due giudici indipendenti,
\texttt{runs/p1\_c2/attempts/packet\_001.json}) non ha trovato passaggi
non giustificati.

\subparagraph{4. Fonti e contributo}\label{fonti-e-contributo}

\begin{itemize}
\tightlist
\item
  Nessuna fonte usata: prova autocontenuta.
\item
  Posizione rispetto allo stato dell'arte: No literature found for this
  lemma specifically (arXiv search: angles between lines / Fejes Toth);
  it is an elementary statement and is proved here from scratch.
\item
  Contributo: la prova qui scritta è del team, per esteso.
\end{itemize}

\subparagraph{5. Limiti e parti
irrisolte}\label{limiti-e-parti-irrisolte}

Nessuna: la parte è dimostrata. Gap dichiarati: nessuno.

\subsubsection{Decisioni degli agenti}
\paragraph{Run \texttt{p1\_c2}}
\begin{description}
\item[Iterazione 1 — attempt\_001]
\textbf{Researcher}: famiglia \texttt{direct\_proof}, sotto-obiettivo: Prove the chain lemma by induction on m, projecting the penultimate vector onto the orthogonal complement of the last one and using the spherical triangle inequality arccos(cos b cos p) \textless{}= b + p.. Stato dichiarato: \texttt{CELL\_SOLVED\_CANDIDATE}.
\emph{Perché questa scelta}: Induction on m is natural because removing the last vector and projecting the penultimate one onto x\_m\textasciicircum{}perp produces a chain one shorter in a space one dimension lower; the loss of angle in this projection is controlled exactly by the right-angle spherical triangle inequality. The degenerate case x\_\{m-1\} = +-x\_m is handled separately by choosing an auxiliary orthogonal unit vector.
\emph{Rispetto allo stato dell'arte}: No literature found for this lemma specifically (arXiv search: angles between lines / Fejes Toth); it is an elementary statement and is proved here from scratch.
\textbf{Referee}: verdetto \texttt{UNKNOWN\_STATUS} (review\_status \texttt{READY\_FOR\_HUMAN}).
\emph{Prossimo passo richiesto}: Human reviews the exact target, proof and evidence, then approves explicit claims
\end{description}

\subsubsection{Codice su cui poggia il tentativo}
Nessun codice nei tentativi.

\section{arXiv literature consulted}
\begin{thebibliography}{99}
\bibitem{1801.07837v1} Dmitriy Bilyk, Ryan W Matzke. \emph{On the Fejes Tóth Problem about the Sum of Angles Between Lines}. arXiv:1801.07837v1 (2018). \url{http://arxiv.org/abs/1801.07837v1}
\bibitem{2007.08698v2} Tongseok Lim, Robert J. McCann. \emph{On Fejes Tóth's conjectured maximizer for the sum of angles between lines}. arXiv:2007.08698v2 (2020). \url{http://arxiv.org/abs/2007.08698v2}
\bibitem{1204.3850v1} Jérémie Chalopin, Shantanu Das, Yann Disser, Matúš Mihalák, Peter Widmayer. \emph{Simple Agents Learn to Find Their Way: An Introduction on Mapping Polygons}. arXiv:1204.3850v1 (2012). \url{http://arxiv.org/abs/1204.3850v1}
\end{thebibliography}
\end{document}
